% 本书原创：直接方法的存在性链条与严格凸性的独立职责。
\begin{figure}[htbp]
\centering
\begin{tikzpicture}[>=stealth,line width=.7pt,
 every node/.style={font=\small},
 directstep/.style={draw=AnalysisBlue,fill=AnalysisBlue!6,rounded corners=2pt,
              align=center,text width=6cm,inner sep=4pt,minimum height=1.05cm}]
\node[directstep] (sequence) at (0,0)
 {极小化序列\\\(u_n\in K,\quad J(u_n)\to m\)};
\node[directstep] (bounded) at (0,-1.5)
 {强制性\\\((u_n)\) 在 \(X\) 中有界};
\node[directstep] (weak) at (0,-3)
 {自反性与弱紧性\\\(u_{n_k}\rightharpoonup u\)};
\node[directstep] (limit) at (0,-4.8)
 {约束弱闭与能量弱下半连续\\\(u\in K\)\\\(J(u)\leq\liminf_kJ(u_{n_k})=m\)};
\node[directstep] (exists) at (0,-6.6)
 {极小元存在\\\(J(u)=m\)};
\node[directstep,text width=3.1cm] (unique) at (5.5,-6.6)
 {另加严格凸性\\极小元唯一};
\draw[->] (sequence)--(bounded);
\draw[->] (bounded)--(weak);
\draw[->] (weak)--(limit);
\draw[->] (limit)--(exists);
\draw[->] (exists)--(unique);
\end{tikzpicture}
\caption[直接方法的存在性步骤与唯一性条件]{直接方法的存在性步骤与唯一性条件。图中 \(X\) 为实自反 Banach 空间，\(K\subset X\) 非空且弱闭，\(J:K\to(-\infty,+\infty]\) 强制、弱下半连续且至少在一点有限，\(m=\inf_KJ\in\mathbb R\)。左列得到极小元；右下分支另要求 \(K\) 凸、\(J\) 在其有限能量点之间严格凸，才得到唯一性。}
\label{b2:fig:14-direct-method}
\end{figure}
